\section{Contested Approaches to the Political Economy of the Rise and Fall of the Unidad Popular}
\label{sec:literature_review}

No single framework commands consensus on the political economy of the rise and fall of Chile's Unidad Popular (UP). Instead, rival explanatory traditions have structured the field, each emerging under unequal political conditions and distinct historical contexts. Transnational networks linking the Pontificia Universidad Católica de Chile (PUC) to the University of Chicago formed the neoliberal school, which later acquired state power through the dictatorship that followed the 1973 coup. Neo-structuralism emerged under the dictatorship as a technocratic opposition to Pinochet's regime—specifically in the aftermath of the 1982 debt crisis—before rising to government power and establishing itself as the mainstream economic paradigm during the post-dictatorship order. Contemporary accounts in the field of Historical Political Economy (HPE) entered the debate much later, bringing powerful tools of causal identification to historical questions. However, they did so within a disciplinary environment already shaped by earlier causal claims taken as given, while perspectives anchored in distributive conflict and structural dependency were seldom considered or actively excluded. Dependency theory occupies a different position altogether. It was the intellectual landscape from which the UP programme itself emerged, and the coup institutionally shattered it through violent repression before it could resolve the theoretical problems it had opened up.

These traditions did not simply coexist. They succeeded one another through political conflict without genuinely opposing each other on the intellectual terrain, despite stark declarations regarding the intellectual bankruptcy of the structuralist-dependency school: 

\begin{quote}
	``...the macroeconomic view of the Unidad Popular failed to recognize that their policies would only be able to generate a burst of economic activity that would be unsustainable in the medium term... Needless to say, this view of the way the economy functioned ignored many of the key principles of traditional economic theory.'' \citep[p.~16--17]{DornbuschEdwards1989}.
\end{quote}

Each identified a part of the problem, yet none integrated external transmission, distributive conflict, and macroeconomic dynamics into a single explanatory sequence. This section critically evaluates these traditions. It reconstructs the historical conditions under which each approach became authoritative, clarifies the mechanisms each privileged, and shows which questions each rendered visible and which it displaced, sanitized, or foreclosed. It begins with classical dependency theory, which posed the problem this paper seeks to resolve more sharply than any approach that followed. 

This fragmented intellectual legacy has resulted in a striking scarcity of rigorous empirical evidence concerning the macroeconomic dynamics of the UP's collapse. While the political and institutional histories of the period are extensively documented, the macroeconomic debate has largely relied on retrospective narratives, analysis of correlation without causality, or ideological posturing rather than methodologically robust quantitative analysis. There is a critical lack of empirical studies employing adequate econometric methods to isolate and measure how allegedly \textit{``populist interventions''} ---such as exchange rate overvaluation, fiscal monetization, and generalized price controls---transmitted distributive conflicts and external shocks into systemic macroeconomic failure. By leaving these macroeconomic transmission mechanisms empirically untested, the literature has remained trapped in theoretical stalemates. This paper seeks to bridge this methodological void by moving beyond ideological retrospection, advancing more rigorous empirical identification to the macroeconomic data of the UP era to trace the anatomy of its economic unraveling.

\subsection{Classical Dependency Theory: Interrupted Programme}
\label{sec:A1_classical_deptheory}

Dependency theory did not approach the Unidad Popular from the outside. It helped constitute the intellectual horizon within which the UP became thinkable. Its formation in Chile brought together two lineages that converged, institutionally and theoretically, in Santiago during the long sixties. From one side came the Monthly Review tradition, beginning with \citet{Baran1957}'s analysis of economic surplus under monopoly capitalism and radicalized by \citet{Frank1967} into the development-of-underdevelopment thesis: underdevelopment was not a residual stage but an active product of incorporation into the world capitalist system. From the other came the structuralist school built within the headquarters of the Economic Commission for Latin America (ECLAC), led by Raúl Prebisch, which had already framed Latin American development through centre--periphery relations, external bottlenecks, and the historical asymmetries of international trade.

At the Centro de Estudios Socioeconómicos (CESO) of the Universidad de Chile, 
these two lineages fused into an active research programme. Following the 1964 military coup 
in Brazil, the arrival of Theotônio dos Santos and Vania Bambirra as political exiles in Santiago 
catalyzed this institutional convergence \citep{CardenasSeabra2022}. Collaborating at CESO 
alongside André Gunder Frank, Orlando Caputo, and Roberto Pizarro—and maintaining close dialogue 
with Ruy Mauro Marini in Concepción—they transformed dependency from a general critique of 
peripheral capitalism into an effort to specify the historical forms linking regional subordination, 
capital accumulation, and political strategy \citep{DosSantos1970,Bambirra1971,CaputoPizarro1970,Kay2010}. 
Far from remaining an academic abstraction, this intellectual collective directly shaped the policy 
architecture of the Chilean left: among Dos Santos and Bambirra's closest collaborators and disciples 
were Sergio Ramos and Alexis Guardia, the primary economic drafters of the Unidad Popular programme 
representing the Communist and Socialist parties, respectively.

The programme's analytical wager was clear. Peripheral capitalism followed a distinct 
developmental dynamic structured by external constraints, a domestic bourgeoisie structurally 
incapable of autonomous national accumulation, and reproduction patterns that continually 
re-entrenched subordination. Within this framework, socialist transformation appeared as an 
objective developmental necessity grounded in the physical limits of peripheral capital formation. 
Yet the tradition never achieved theoretical unity. Marini's theory of superexploitation pushed 
the analysis toward a world-market logic where the systematic compression of labor below the value 
of its reproduction conditioned domestic capital accumulation \citep{Marini1973}. Conversely, 
\citet{CardosoFaletto1979} insisted that dependency took historically contingent forms governed 
by domestic class alliances and the relative autonomy of the capitalist state. This tension proved 
decisive. At stake was whether domestic political mediation constituted an independent object 
of analysis or simply mediated world-market imperatives. Classical dependency posed the question 
with great force but never resolved it.


That unresolved tension did not prevent the tradition from producing the richest 
contemporaneous political-economy accounts of the Unidad Popular. \citet{DeVylder1976}'s \textit{Allende's Chile} 
provided the seminal participant-observer account, correctly locating the material foundation 
of the external constraint in Chile's extreme physical import dependency: over $90\%$ of capital 
equipment was imported, and more than $95\%$ of machinery in large-scale copper mining originated 
in the United States. Yet De Vylder's framework suffered from an important analytical compromise. 
Confronted with orthodox arguments attributing the crisis entirely to domestic mismanagement, 
De Vylder retreated into an internal-blame explanation, asserting that: 
\begin{quote}
	\emph{``the problem was the  insatiable imports, not the US blockade nor the lack of credit''} \citep[p.~106]{DeVylder1976}. 
\end{quote}

To sustain this concession, De Vylder pointed to gross nominal credit lines extended by the Soviet 
Union and Eastern Europe ($>\$400$~million) as ostensible proof that external financing remained 
abundant, conflating gross nominal commitments with usable, convertible liquidity.

\citet{Stallings1978}  overturned De Vylder's concession by uncovering 
the \textit{qualitative non-fungibility of international credit}. Drawing on primary Central 
Bank foreign exchange ledgers, she demonstrated that bilateral socialist credits were 
non-convertible tied lines earmarked for specific industrial equipment. They could neither 
purchase patented American spare parts for Kennecott copper haul trucks nor settle immediate 
spot cash contracts for wheat and petroleum on Western commodity markets. Concurrently, private 
U.S. commercial banks slashed revolving short-term trade credit lines from $\$220$~million to 
$\$35$~million, while the Export-Import Bank froze disbursements. These debates around balance-of-payments constraints did not unfold in isolation. They formed part of a 
vibrant contemporaneous intellectual landscape captivated by the \textit{vía chilena al socialismo}---the 
historical attempt to achieve revolutionary socialist transformation through 
bourgeois democratic institutions. Across this field emerged a rich empirical corpus tracking 
the crisis from complementary angles: distributional conflict across successive administrations, 
worker and peasant self-organization from below, the logistical spare-parts embargo, the monetary 
and balance-of-payments mechanisms through which external pressure became domestic crisis, and 
the strategic dilemmas of the state under conditions of emerging dual power 
\citep{StallingsZimbalist1975,EspinosaZimbalist1978,Kay1978,PetrasMorley1975,GriffithJones1981,Sideri1979,SweezyMagdoff1974,ZavaletaMercado1974,Winn1986}.  The enduring methodological strength of classical dependency lay precisely here: it refused to  treat external imperial intervention and domestic class struggle as competing explanations,  approaching them instead as interacting, co-constitutive dimensions of the same historical process.

That achievement also marked the analytical limit dependency theory could not overcome 
in its classical form. Confined to a bilateral nation-state framework (the U.S. versus Chile), it omitted 
the international monetary architecture of Bretton Woods and global stagflation, while lacking a formal 
macroeconomic specification linking external balance-sheet shocks to domestic distributive conflict and 
capacity utilization.  Yet the defeat of dependency theory was twofold. Politically, the coup dismantled the institutions, networks, and research centres in which the programme had been elaborated. Intellectually, the 
interruption arrived before the tradition could work through its own internal tensions. No superior 
explanatory framework refuted it in a scientific terrain. It broke under conditions in which theoretical incompletion coincided with the collapse of Chile's liberal democracy, the failure of democratic socialism to 
resolve the crisis under organized domestic and international sabotage, and the installation of a 
subsequent civic-military dictatorship. Subsequent approaches entered a field already reorganized 
by force, not reason. The next school to become authoritative in Chilean political economy was not 
the one that best resolved the problem dependency had identified, but the one that arrived coordinated 
with the counterrevolutionary historical bloc that constituted the reconstruction of the state and civil society \citep{Casals2021}.

\subsection{The Rise of the Chicago Boys and the Neo-liberal School}
\label{sec:A2_neoliberal_chicago}

The neoliberal school did not emerge from within the intellectual world of the Unidad 
Popular as one competing interpretation among others. It entered through a different route: as a 
transnationally trained programme of economic reconstruction that acquired authority through the 
counterrevolutionary reorganization of the Chilean state. Its genealogy is well established. The 1956 
agreement between the Pontificia Universidad Católica de Chile and the University of Chicago opened 
the institutional channel through which Milton Friedman, Arnold Harberger, and Gary Becker trained 
a cohort of Chilean economists in monetarism, public finance orthodoxy, and neoclassical price theory. 
By the late 1960s, this network had evolved from an academic outpost into a politically embedded 
project, closely aligned with business associations and the Chilean right, oriented toward the systemic 
restructuring of Chilean capitalism \citep{Valdes1995,Fischer2009,Clark2017}.

That programme found its canonical expression in \textit{El Ladrillo}\footnote{\textit{El Ladrillo} (``The Brick'') was drafted clandestinely prior to the September 1973 coup by Chicago-trained 
economists coordinated by Emilio Sanfuentes, with contributions from Sergio de Castro and Pablo Baraona, 
in liaison with naval officers and the leadership of the private employers' association \textit{Sociedad de Fomento Fabril} (SOFOFA)  \citep{DeCastro1992}.}. The text codified in advance the macroeconomic blueprint on which the post-UP order would be constructed. It operated as a counterrevolutionary vanguard design: an 
operational policy manual prepared prior to the coup to direct the reconstruction of Chilean capitalism 
the moment the popular government was overthrown.

Its diagnosis identified a familiar catalogue of structural problems: chronic inflation, 
recurrent external deficits, agricultural stagnation, low productivity growth, and distortions 
associated with state-led industrialization. Several of these observations captured real tensions within 
Chilean accumulation, and they merit serious engagement. The argument that prolonged relative-price 
controls hindered export diversification and exacerbated balance-of-payments vulnerability cannot 
simply be dismissed, nor can the recognition that import-substitution industrialization had reached severe 
structural exhaustion \citep{DeCastro1992}. Yet the broader theoretical hypothesis holding these observations 
together was much weaker: the claim that the crisis of the UP merely culminated a secular drift toward the 
excessive socialization of the economy. That thesis recoded a complex conjunctural crisis into a normative 
indictment, naturalizing market discipline rather than explaining the empirical dynamics of the breakdown.

Here lay the neoliberal school's structural limitation. It registered macroeconomic imbalances 
only by systematically stripping them of their political-economy. Functional income distribution 
appeared not as structural class conflict, but as a technical distortion of price signals. The external 
crisis was treated as the endogenous consequence of fiscal and monetary mismanagement, never as the relational 
interaction between domestic class struggle and international financial subordination and destabilization. The October 1972  employers' lockout, the private commercial credit blockade, the spare-parts embargo, and covert foreign destabilization could not enter the account without undermining the framework's own core hypothesis and axioms. Neoliberalism did not omit these factors by accident. Its explanatory logic required them to disappear, since acknowledging them would have undermined the premises that the crisis proved the inherent impossibility of socialist transformation.

The neoliberal interpretation became authoritative under historical conditions that were not 
primarily intellectual. Its ascendancy reflected no superior theoretical integration of external transmission, 
distributive conflict, and macroeconomic dynamics. A programme formulated before the coup entered a terrain 
reorganized by military force, and the dictatorship installed it as the common sense of national economic 
reconstruction. That restructuring violently dismantled collective bargaining, political mediation, and 
labor organizations while relying, paradoxically, on the centralized administrative apparatus of the developmental  state it claimed to dismantle \citep{Clark2017,Bockman2019}. The neoliberal school was less a scientific solution  to the developmental dilemmas of the periphery than a victorious political project presenting ruling class backlash as economic science. Its enduring legacy lies precisely here: it converted a conjunctural crisis resolved by state terror into a universalized economic doctrine asserting the impossibility of any non-market alternative \citep{Espinosa2021}.

\subsection{The Neo-Structuralist School: Selective Sanitization}
\label{sec:A3_neo_structuralism}

The neo-structuralist school was an intellectual accomplishment of the post-dictatorship 
order rather than a simple continuation of classical structuralism. It established its institutional 
base during the dictatorship, primarily within the CIEPLAN think tank, where former structuralists and 
democratic opposition economists critically monitored the Chicago Boys' reforms while formulating an 
alternative programme for the democratic transition. That genealogy explains much of the school's 
subsequent authority. The CIEPLAN economists were more than political moderates who had tempered their 
commitments under authoritarian rule \citep{Silva1991}. They constituted the modernizing technocratic 
fraction through which the post-authoritarian centre-left could challenge orthodox monetarism while 
operating strictly within market-disciplined, export-oriented peripheral capitalism. 
Neo-structuralism entered the field as a politically functional renovation calibrated to manage Chile's 
subordinate insertion into the globalized economy \citep{Sklair2000,Leiva2008,Maillet2016}.

The central analytical move of that renovation was the abandonment of the classical 
surplus approach. Where Prebisch and the foundational ECLAC tradition had centered the production, 
appropriation, and reinvestment of the economic surplus under peripheral conditions, neo-structuralism 
recast underdevelopment into the mainstream idiom of market imperfections, institutional bottlenecks, 
and pragmatic state regulation. This shift did not merely modernize structuralism; it altered its 
analytical object. The external constraint remained, but strictly as a macroeconomic financing ceiling 
rather than as a structural relation of surplus extraction and dependent reproduction. By reducing a 
theory of historically structured subordination to a technical problem of policy calibration, the 
school depoliticized the external barrier. As \citet{Leiva2008} demonstrates, without the surplus 
approach, the external constraint survives only in a thinned accounting form: legible as a foreign 
exchange deficit, but stripped of the class relations that generate and sustain it. Neo-structuralism 
did not forfeit empirical rigor; it forfeited the conceptual framework required to theorize the 
balance of payments and distributive conflict as unified moments of the same accumulation process.

The school's account of the UP crystallized most authoritatively in the work of Sergio Bitar. 
Originally published in 1979, \citet{Bitar2013} provided the most sophisticated neo-structuralist 
reconstruction of the popular government. Its principal merit lay in rejecting an exclusively monetary 
explanation of the external accounts. Bitar sought to integrate foreign exchange scarcity, political 
opposition, and institutional transformation into a single dynamic narrative. His formulation of a 
``viability area'' remains a genuine analytical contribution, demarcating the narrow economic margin 
within which the UP programme could have advanced without precipitating macroeconomic collapse. Yet 
that contribution also established the school's defining silence. Across multiple editions, Bitar never 
substantively engaged the contemporaneous dependency literature that had developed across the same 
historical terrain. Foundational accounts by \citet{DeVylder1976} and \citet{Stallings1978} were already 
widely available, yet Bitar avoided confronting them. This silence was neither a bibliographic omission 
nor a conspiracy. It reflected the strategic selectivities accompanying the abandonment of the surplus 
approach \citep{Leiva2008}, marking the point at which dependency's central questions were cordoned 
off from post-dictatorship economic discourse.

That theoretical cordon became hegemonic through the populist-macroeconomic paradigm. 
\citet{LarrainMeller1991} codified the UP into a stylized four-phase trajectory of initial reactivation, 
emerging bottlenecks, severe destabilization, and ultimate collapse. \citet{DornbuschEdwards1991} then 
generalized this sequence into a regional archetype of ``macroeconomic populism.'' This paradigm achieved 
wide influence by providing a clear narrative, identifiable turning points, and cautionary lessons that 
circulated throughout Latin America during its democratic transitions. Yet it operated through a 
definitional foreclosure: by classifying redistributive policies categorically as \textit{``macroeconomic 
populism,''} the framework ruled out the possibility that wage-led growth could ever be structurally 
sustainable under peripheral conditions. \citet{Meller2000} later cemented this interpretation as 
the common sense of the UP: an unforced tragedy driven by policy mismanagement and political polarization. 
The result was a reading self-critical in tone but exculpatory in effect, reducing the crisis to domestic 
excess while treating external financial blockades and domestic logistical sabotage as secondary noise. 
The contrast between Sergio \citet{RamosSergio1979} and Joseph \citet{RamosJoseph1980} illustrates this epistemological 
divide: the former treated inflation as a structural contradiction internal to a socialist transition 
under acute class struggle, whereas the latter reduced the crisis to an engineering problem of disequilibrium 
management and stabilization failure.

Neo-structuralism thus substituted a technocratic problem of governance for the 
conflict-centered political economy generated during the UP itself \citep{LarrainMeller1991,DornbuschEdwards1991,Meller2000}. 
The school's analytical limits were not accidental blindspots; they were strategic exclusions built 
into the conditions of its post-authoritarian legitimacy. Its primary historical function was to 
preserve the structuralist language of external constraints while adapting it to the macroeconomic 
governance of a neoliberalized economy, sustaining two decades of centre-left political stability. In 
doing so, it transformed the thousand days of Allende's government into a moral tale of redistributive 
overreach. Neo-structuralism became the progressive framework through which the post-dictatorship order  could be managed and stabilized without reviving the class-conflict categories that made earlier structuralist and dependency traditions dangerous to ruling-class hegemony \citep{Leiva2008}.

This critique does not dismiss the empirical contributions of scholars such as Bitar or 
Meller; it specifies the political-economic compromises and strategic selectivities that resulted from their synthesis. The coordination of monetary, structural, and political variables into a coherent narrative remain important contributions to the debate. However, the systematic exclusion of class struggle for the economic surplus was the price of its institutionalization during the democratic transition and the post-dictatorship period. That compromise defines neo-structuralism's enduring legacy. It established the default macroeconomic memory inherited by subsequent scholarship, leaving historical political economy equipped with sophisticated empirical tools, but lacking a coherent theory of the structural mechanisms at work.

\subsection{Historical Political Economy: Methods Under Inherited Macro Theory}
\label{sec:A4_historical_PE}

Historical Political Economy (HPE) represents a genuine methodological advance in the study 
of economic history. By applying the credibility revolution's identification toolkit to historical inquiry, 
it has introduced a level of empirical precision rarely achieved in earlier debates. Its institutionalization 
around the \textit{Journal of Historical Political Economy} and the \textit{Oxford Handbook of Historical 
	Political Economy} reflects a broader effort within economics to reclaim historical questions without sacrificing 
rigorous causal identification \citep{Gailmard2021,JenkinsRubin2024,CharnyshFinkelGehlbach2023}. Applied to twentieth-century 
Chile, this agenda has generated an impressive corpus examining price controls, real wages, land redistribution, 
and political violence.

Yet HPE did not generate the archival terrain on which it operates. That foundation was 
constructed by a parallel renewal in Chilean economic history, which has reconstructed long-run time series 
on prices, trade, living standards, and distribution. The archival work of Mario Matus, Manuel Llorca-Jaña, 
Claudio Robles, and their collaborators established the quantitative records that made micro-econometric 
identification feasible. Within this archival foundation, \citet{RodriguezWeber2018} is foundational for the 
empirical architecture of \citet{Polanco2026}: the long-run historical accounts of inequality compiled there 
directly underpin the functional wage-share series ($\omega_t$) harmonized by \citet{Astorga2023} and utilized 
throughout Chapters~2 and~3. By situating the Unidad Popular within a secular century-long cycle of distribution 
and explicitly demonstrating how political coercion shaped income shares, Rodríguez Weber provides a central 
historical bridge between institutional economic history and structural macroeconomic questions.

In contrast to descriptive economic history, HPE proper focuses on estimating discrete causal 
effects. \citet{AldunateGonzalezPrem2024} exploit firm-level records to identify the credit-allocation 
shocks induced by U.S. financial warfare, while \citet{JaimovichToledo2021} trace contemporary Mapuche territorial 
conflict to the localized geography of agrarian expropriation and military counter-reform. \citet{GirardiBowles2018} 
provide a distinctive contribution: their high-frequency event study of the Santiago stock exchange demonstrates 
that financial wealth-holders responded systematically to shifts in the political threat to private property, 
corroborating classical political economy's focus on capitalist investment strikes. These studies represent 
sophisticated empirical contributions. However, they perform fundamentally different analytical work: some 
build archival ledgers, others isolate localized causal effects, while few---if not \citet{GirardiBowles2018} 
alone---examine macroeconomics and institutions through a theoretical framework centered on the structural 
power of capital over the economy.

This analytical demarcation is most evident in the work of Felipe González and his collaborators, 
who define the contemporary frontier of Chilean HPE \citep{Gonzalez2013,GonzalezVial2021,AldunateGonzalezPrem2024,GonzalezPremUrzua2020,BautistaEtAl2023,GonzalezPrem2026Nat,GonzalezPrem2026Trans}. 
Their research is methodologically exemplary. They have unearthed primary archival records, formulated transparent 
quasi-experimental designs, and resolved contested empirical questions surrounding land reform, industrial 
requisitions, covert foreign intervention, and the institutional legacies of the military dictatorship. Yet this 
corpus also marks the field's primary constraint: the macroeconomy appears primarily as a passive background 
upon which institutional interventions and policy shocks are measured, rather than as an organic, contradictory 
system of capital accumulation, surplus extraction, and crisis. Even where the archival evidence touches systemic 
macroeconomic phenomena, the theoretical interpretation is routinely imported from post-dictatorship mainstream 
narratives. The result is scholarship of exceptional empirical precision that displaces political economy in the 
classical sense into an applied microeconomics of institutional interventions.

This limitation reflects the field's disciplinary formation. HPE arrived at the Chilean case after 
the historical clash between the dependency, neo-structuralist, and neoliberal paradigms had already been settled 
by state coercion. It inherited the depoliticized macroeconomic vocabulary of the 
post-dictatorship consensus, which reduced the UP's crisis to a cautionary tale of macroeconomic populism and 
unforced balance-of-payments mismanagement. The research programme developed in this paper builds upon HPE's 
archival and empirical rigor, while restoring what applied microeconomics omits: an integrated theoretical framework  capable of modeling external balance-of-payments vulnerability, distributive conflict, and domestic monetary  accommodation as mutually conditioning moments of peripheral capital accumulation.

\subsection{The Revival of Dependency Theory as a Research Programme}
\label{sec:A5_rev_deptheory}

The return to dependency theory is not an exercise in intellectual nostalgia. 
It recovers a research programme whose historical retreat signaled a profound collapse of the region's 
sociological imagination. As \citet{Palma2016} argues, Latin America's critical thinking swung like 
a pendulum from one type of absolute certainty to another: from the mechanical stagnationism of early 
dependency to the uncritical market fatalism of the post-dictatorship consensus. In that shift, the 
region's intellectual landscape was stripped of its capacity to conceive development as a problem of 
collective agency, structural transformation, and sovereign autonomy. Development was degraded from a 
political struggle over the production, appropriation, and reinvestment of the economic surplus into an 
administrative exercise in macroeconomic compliance. The analytical task today is neither to revive the 
rigid certainties of the 1970s nor to surrender to technocratic fatalism, but to reconstruct dependency 
as an empirically generative, progressive framework.

\citet{Kvangraven2021} provides the epistemological architecture for this reconstruction 
by formalizing dependency theory through Imre Lakatos's methodology of scientific research programmes. 
Kvangraven refutes the mainstream caricature of dependency as a monolithic dogma by demonstrating that 
the tradition is structured around a coherent, shared \textit{hard core}: the premise that global capitalist 
accumulation systematically reproduces hierarchical structures between core and periphery that constrain 
the developmental space of subordinate social formations. The scientific vitality of the programme does not depend on dogmatic consensus regarding specific historical models; rather, it hinges on whether its \textit{protective belt of auxiliary hypotheses} can generate \textit{progressive problemshifts}---predicting novel empirical facts and specifying concrete causal mechanisms---rather than degenerating into defensive ad-hoc immunizations. By formalizing this distinction, she clarifies how dependency can assimilate new historical evidence within the context of the concrete realities of global capitalism without abandoning its core ontological commitments.

\citet{MadariagaPalestini2021} sharpen this epistemological reconstruction by distinguishing 
between \textit{situations of dependency} and \textit{mechanisms of dependency}: i.e., between historically specific  descriptive configurations and the causal processes that reproduce external constraints. This 
distinction has energized a contemporary renewal in Santiago, where scholars have operationalized dependency 
theory to analyze rentier accumulation, export stagnation, and dependent financialization in present-day Chile 
\citep{Ahumada2019,Ruiz2021,AhumadaTorres2024,Sossdorf2024}. Within this revival, recovering the historical memory  of the Unidad Popular is not an exercise in political nostalgia, but an inquiry into the structural conditions 
under which sovereign control over the surplus and democratic self-determination were attempted and violently 
suppressed \citep{Arboleda2023,AhumadaChang2025}. Yet for the macroeconomic analysis of the UP itself, the 
central task remains unfinished: while classical accounts identified the concrete \textit{situation} of Chilean 
vulnerability (import dependency and bilateral credit retaliation), they lacked a formal \textit{mechanism} 
linking peripheral balance-of-payments asphyxiation to the structural crisis of the core's international monetary order.

This missing transmission mechanism can be formulated by synthesizing the political economy 
of global dollar hegemony with the institutional mechanics of the dissolution of Bretton 
Woods. \citet{Vernengo2021} provides the structural foundation for why Bretton Woods collapsed: a macro-political drive by U.S. elites to escape the discipline of gold convertibility, discipline domestic labor through international capital mobility, and establish an unconstrained fiat dollar standard. In turn, \citet{Zoeller2019} demonstrates how and when this rupture unfolded: rather than a predetermined grand design, closing the gold window was a contingent executive maneuver forced by IMF institutional veto points and European surplus resistance. The core executed an emergency exit on August 15, 1971, rapidly consolidating a floating dollar regime that eliminated its own reserve discipline while transmitting an external shock to peripheral economies.

For Chile, this historical timing was decisive. Nixon's unilateral closure of the gold window 
occurred exactly three days after the U.S. Export-Import Bank formally denied credit lines to the Unidad Popular on August 12, 1971. Peripheral accumulation and the \textit{``via Chilena al Socialismo''} was thus caught in converging external constraints: the bilateral credit non-fungibility documented by \citet{Stallings1978} colliding with the systemic dissolution of the global monetary regime. Incorporating this dual transmission into the protective belt of dependency theory resolves the theoretical ambiguities that stalled its classical formulation. It equips the tradition with an exogenous historical event (Nixon's unilateral suspension of gold convertibility) as a macroeconomic mechanism capable of explaining how external balance-sheet shocks triggered domestic supply paralysis and forced monetary accommodation, directly setting up the comparative evaluation of competing hypotheses developed in this paper.  

\subsection{Epistemological Critique: The Illusion of Causality, Narrative License, and Collider Traps}
\label{sec:epistemological_critique}

The persistent divergence among these explanatory traditions cannot be resolved through historiographical narrative alone. As \citet{BlancoMatute2018} establish, human causal reasoning is systematically vulnerable to the \emph{illusion of causality}: a robust cognitive bias wherein observers detect strong causal links between co-occurring events even when the statistical contingency is zero or driven by unobserved common causes. In the economic historiography of the Unidad Popular, orthodox monetarist and Austrian accounts observe a conspicuous historical coincidence---the election of a socialist government in late 1970, expansionary fiscal-wage policy in 1971, and the explosion of inflation in 1972--1973. Relying on intuitive heuristics, they assign exclusive causal responsibility to domestic monetary issuance ($g_{M,t} \to \pi_t$), treating international monetary hierarchy, terms-of-trade shifts, and geopolitical destabilization as irrelevant background noise \citep{Edwards2024,EspinosaVejar2025}.

This causal illusion is reinforced by what \citet{IschEtAl2026} identify as \emph{narrative license} and \emph{overreaching causal claims} in historical social science. Orthodox accounts dress up doctrinal priors as ironclad empirical laws without testing alternative causal mechanisms nor specifying counterfactual identification paths. For example, \citet{EspinosaVejar2025} assert an Austrian Business Cycle mechanism wherein artificial credit expansion generated an unsustainable capital malinvestment boom. Yet they advance this claim without presenting positive evidence of gross fixed capital deepening ($\Delta K$), ignoring that Chilean fixed machinery investment contracted during 1971 while capacity utilization reached historical peaks. Narrative license allows the theorist to impute unobserved microeconomic mechanisms (capital lengthening) to explain macroeconomic outcomes while dismissing external interventions as quantitatively negligible.

Methodologically, orthodox accounts fall into what causal inference defines as a \emph{collider trap} \citep{Pearl2009}. A collider trap takes place when analysts focus on a symptom that is actually the shared outcome of two separate causes. By isolating their attention on this shared outcome, they distort the relationship between the causes, often falsely attributing the symptom entirely to just one of them. In narratives of the UP crisis, this shared outcome---the collider---is the widespread presence of retail shortages and food queues ($C$) \citep{Edwards2024}. 

Orthodox accounts treat these empty shelves as definitive proof of \textit{``repressed inflation,''} attributing them solely to the government printing money to stimulate demand ($U_1$). However, these shortages were simultaneously driven by a second, independent factor: the systematic hoarding and lockouts organized by private wholesale distributors and trucking strikes ($U_2$). Both state demand stimulus ($U_1$) and private supply lockouts ($U_2$) collided to produce empty retail shelves ($C$):
\begin{equation}
	U_1 \longrightarrow [C] \longleftarrow U_2 \quad \Longrightarrow \quad U_1 \not\perp\!\!\!\perp U_2 \mid C
\end{equation}
By treating retail shortages as a neutral, objective measure of excess liquidity (conditioning on $C$), orthodox accounts commit a collider error of causal identification. This methodological blind spot produces a spurious conclusion, allowing them to blame public deficit monetization for a supply paralysis that was heavily driven by the political actions of private business federations.

To transcend these causal illusions and collider traps, I adapt the methodological approach of \citet{KingSznajder2006}, who formally tested competing economic theories to explain Poland's growth trajectory during neoliberalization. In this paper, I treat Monetarism ($H_1$), the Austrian Business Cycle ($H_2$), and Dependency Theory ($H_3$) as rival scientific research programmes. Each possesses a distinct theoretical core, auxiliary assumptions, and falsifiable empirical predictions. I submit the central causal claims of these competing traditions to a systematic comparative evaluation across multi-frequency time series data.
